\section{Conclusion}
\label{sec:conclusion}

We set out to test whether language models exhibit \emph{graded} hidden robustness to
adversarial fine-tuning --- whether per-prompt behavioral drift decreases with a prompt's
semantic distance from the fine-tuning domain, forming low-drift basins on orthogonal tasks.
We built \method, a checkpoint-level trajectory-divergence diagnostic, fine-tuned a real
instruct model on the canonical \insecure-code dataset against a matched benign \secure
control, and measured per-prompt drift across a graded probe set.

\para{Key findings.} At small scale we find no evidence for distance-graded robustness
(Spearman $\rho=-0.08$, $p=0.78$) and no orthogonal-task basin --- in fact orthogonal-capability
prompts drift the \emph{most} (0.553) and a factual code-safety pocket drifts the \emph{least}
(0.045). Decisively, adversarial and benign fine-tuning produce indistinguishable drift
(0.356 vs.\ 0.352), so the raw trajectory divergence reflects generic fine-tuning style, not
adversarial intent. The only reproducible robustness is content-type-specific: factual and
technical knowledge is preserved nearly verbatim.

\para{Takeaway.} Trajectory-divergence metrics are diagnostic of \emph{adversarial}
robustness only when differenced against a matched benign-fine-tuning control; raw drift
overstates misalignment, and orthogonality to the fine-tuning objective does not confer
behavioral robustness.

\para{Future work.} With adequate compute, the natural next steps are to (1) scale to models
where emergent misalignment is known to appear (\eg Qwen2.5-0.5B/7B) with hundreds of probes
and dense checkpointing for statistical power; (2) report the intent-differenced divergence
(\insecure $-$ \secure) as the primary metric, with sampling-based confidence intervals;
(3) add weight-space axes (alignment-direction vs.\ orthogonal displacement) to test whether
behavioral robustness maps to weight-space basins; and (4) test the instability prediction
of \citet{springer2026geometry} directly by checking whether any early-training robustness
collapses later in fine-tuning.
